%% kapa: exam=1 round=31 year=2020 period=1 form=A subject=civil_law subject_ko=민법 seq=8 number=8
%% source: Q-Net 감정평가사 기출문제 (cst00309) artlSeq=5207219 file=extracted/1008/5207219/A형/A형/2020년도 제31회 감정평가사 1차 시험_1교시_A형.pdf
%% conversion: kordoc 4.24.51 -> tools/gen1.py (auto); review=proofread-sonnet; flags=none
\begin{question}{8}
강행규정에 위반되어 그 효력이 인정되지 않는 것을 모두 고른 것은? (다툼이 있으면 판례에 따름)
\begin{framedtext}
ㄱ. 제3자가 타인의 동의를 받지 않고 타인을 보험계약자 및 피보험자로 하여 체결한 생명보험계약 \par
ㄴ. 건물의 임차인이 비용을 지출하여 개조한 부분에 대한 원상회복의무를 면하는 대신 그 개조비용의 상환청구권을 포기하기로 하는 약정 \par
ㄷ. 사단법인의 사원의 지위를 양도․상속할 수 있다는 규약 \par
ㄹ. 승소를 시켜주면 소송물의 일부를 양도하겠다는 민사소송의 당사자와 변호사 아닌 자 사이의 약정 \par
\end{framedtext}
\begin{choices}
  \item ㄱ, ㄴ
  \item ㄱ, ㄷ
  \item ㄱ, ㄹ
  \item ㄴ, ㄷ
  \item ㄷ, ㄹ
\end{choices}
\end{question}
\answer{3}
